% Лекц 6: семинар ба гэрийн даалгавар. Бодлогууд: Г.Н. Берман, "Сборник задач по курсу
% математического анализа" (22-р хэвлэл). [Б.N] нь номын бодлогын дугаар.

\ProblemSet{Семинар 6}
\begin{frame}[allowframebreaks]{Семинар 6}
\pgroup{Хязгаарын тодорхойлолт ($\varepsilon$--$\delta$ хэлээр).}
\begin{plist}
\prob{190}{$y=x^2$. $x\to2$ үед $y\to4$. $|x-2|<\delta$-аас $|y-4|<\varepsilon=0{,}001$ мөрдөхийн тулд $\delta$ ямар байх вэ?}
\prob{192}{$y=\frac{x-1}{2(x+1)}$. $x\to3$ үед $y\to\frac14$. $|x-3|<\delta$-аас $\bigl|\frac14-y\bigr|<0{,}01$ мөрдөх $\delta$-г ол.}
\prob{193}{$x\to\frac{\pi}{2}$ үед $\sin x\to1$ болохыг батал. $1-\sin x<0{,}01$ байхын тулд $x$ нь $\frac{\pi}{2}$ цэгийн орчимд ямар нөхцөл хангах ёстой вэ?}
\prob{194}{$\displaystyle\lim_{x\to\infty}\frac{1}{x^2+1}=0$. $|x|>N$-ээс $y<\varepsilon$ мөрдөхийн тулд $N$ ямар байх вэ?}
\prob{195}{$x\to+\infty$ үед $y=\frac{x^2-1}{x^2+3}\to1$. $|x|>N$-ээс $|y-1|<\varepsilon$ мөрдөх $N$-г ол.}
\prob{213}{$x\to0$ үед $y=\frac{x}{x+1}\to0$ болохыг батал. $|y|<10^{-4}$ байхын тулд $x$ ямар нөхцөл хангах вэ?}
\prob{214}{$x\to+\infty$ үед $y=\sqrt{x+1}-\sqrt x\to0$ болохыг харуул. $x>N$ үед $y<\varepsilon$ байхын тулд $N$ ямар байх вэ?}
\end{plist}
\pgroup{Төгсгөлгүй их функц.}
\begin{plist}
\prob{198}{$x\to0$ үед $y=\frac{1+2x}{x}\to\infty$. $|y|>10^4$ байхын тулд $x$ ямар нөхцөл хангах вэ?}
\prob{199}{$x\to3$ үед $y=\frac{x}{x-3}$ төгсгөлгүй их болохыг батал. $|y|>1000$ байхын тулд $x$ ямар байх вэ?}
\prob{200}{$x\to1$ үед $y=\frac{1}{(x-1)^2}$ хязгааргүй өснө. $|x-1|<\delta$-аас $\frac{1}{(x-1)^2}>N=10^4$ мөрдөх $\delta$-г ол.}
\end{plist}
\pgroup{Хязгаарыг ол.}
\begin{plist}
\prob{268}{$\displaystyle\lim_{x\to2}\frac{x^2+5}{x^2-3}$}
\prob{272}{$\displaystyle\lim_{x\to1}\frac{x^2-2x+1}{x^3-x}$}
\prob{274}{$\displaystyle\lim_{x\to1}\frac{(x-1)\sqrt{2-x}}{x^2-1}$}
\prob{275}{$\displaystyle\lim_{x\to\frac12}\frac{8x^3-1}{6x^2-5x+1}$}
\prob{276}{$\displaystyle\lim_{x\to1}\frac{x^3+x-2}{x^3-x^2-x+1}$}
\prob{277}{$\displaystyle\lim_{x\to1}\Bigl(\frac{1}{1-x}-\frac{3}{1-x^3}\Bigr)$}
\prob{281}{$\displaystyle\lim_{x\to\infty}\frac{x^3+x}{x^4-3x^2+1}$}
\prob{283}{$\displaystyle\lim_{x\to\infty}\frac{x^2-1}{2x^2+1}$}
\prob{293}{$\displaystyle\lim_{x\to0}\frac{\sqrt{1+x^2}-1}{x}$}
\prob{296}{$\displaystyle\lim_{x\to5}\frac{\sqrt{x-1}-2}{x-5}$}
\end{plist}
\end{frame}

\ProblemSet{Гэрийн даалгавар 6}
\begin{frame}[allowframebreaks]{Гэрийн даалгавар 6}
\pgroup{Хязгаарын тодорхойлолт, төгсгөлгүй их ба зааглагдсан функц.}
\begin{plist}
\prob{191}{$y=\frac{x^2-1}{x^2+1}$. $x\to2$ үед $y\to\frac35$. $|x-2|<\delta$-аас $\bigl|y-\frac35\bigr|<0{,}1$ мөрдөх $\delta$-г ол.}
\prob{201}{$x\to0$ үед $y=\frac{1}{2^x-1}$ төгсгөлгүй их. $|y|>100$ байхын тулд $x$ ямар тэнцэтгэл биш хангах вэ?}
\prob{202}{$x\to+\infty$ үед $y=\lg x\to+\infty$. $x>M$-ээс $y>N=100$ мөрдөхийн тулд $M$ ямар байх вэ?}
\prob{204}{$y=\frac{x^2}{1+x^4}$ функц бүх тоон шулуун дээр зааглагдсан болохыг батал.}
\prob{205}{$y=\frac{x^2}{1+x^5}$ функц бүх тоон шулуун дээр зааглагдсан уу? $(0,+\infty)$ завсарт зааглагдсан уу?}
\prob{210}{Төгсгөлгүй их функц мөн үү: 1) $f(x)=x\arctg x$, \ $x\to\infty$; \ 2) $f(x)=(2+\sin x)\lg x$, \ $x\to+\infty$; \ 3) $f(x)=(1+\sin x)\lg x$, \ $x\to+\infty$?}
\prob{215}{$x\to\infty$ үед $f(x)$-ийн хязгаар $a$ бол $f(x)=a+\varphi(x)$ ($\varphi$ нь $x\to\infty$ үед төгсгөлгүй бага) хэлбэрт бичигдэхийг батал. $y=\frac{x^3}{x^3-1}$, \ $y=\frac{x^2}{2x^2+1}$, \ $y=\frac{1-x^2}{1+x^2}$-г ийм хэлбэрт бич.}
\end{plist}
\pgroup{Хязгаарыг ол.}
\begin{plist}
\prob{269}{$\displaystyle\lim_{x\to0}\Bigl(\frac{x^3-3x+1}{x-4}+1\Bigr)$}
\prob{270}{$\displaystyle\lim_{x\to1}\frac{x}{1-x}$}
\prob{271}{$\displaystyle\lim_{x\to\sqrt3}\frac{x^2-3}{x^4+x^2+1}$}
\prob{273}{$\displaystyle\lim_{x\to-2}\frac{x^3+3x^2+2x}{x^2-x-6}$}
\prob{278}{$\displaystyle\lim_{x\to2}\Bigl[\frac{1}{x(x-2)^2}-\frac{1}{x^2-3x+2}\Bigr]$}
\prob{279}{$\displaystyle\lim_{x\to1}\Bigl[\frac{x+2}{x^2-5x+4}+\frac{x-4}{3(x^2-3x+2)}\Bigr]$}
\prob{280}{$\displaystyle\lim_{x\to1}\frac{x^m-1}{x^n-1}$ \ ($m,n$ бүхэл)}
\prob{282}{$\displaystyle\lim_{x\to\infty}\frac{x^4-5x}{x^2-3x+1}$}
\prob{284}{$\displaystyle\lim_{x\to\infty}\frac{1+x-3x^3}{1+x^2+3x^3}$}
\prob{285}{$\displaystyle\lim_{x\to\infty}\Bigl(\frac{x^3}{x^2+1}-x\Bigr)$}
\prob{286}{$\displaystyle\lim_{x\to\infty}\Bigl(\frac{x^3}{2x^2-1}-\frac{x^2}{2x+1}\Bigr)$}
\prob{287}{$\displaystyle\lim_{x\to\infty}\Bigl[\frac{3x^2}{2x+1}-\frac{(2x-1)(3x^2+x+2)}{4x^2}\Bigr]$}
\prob{294}{$\displaystyle\lim_{x\to0}\frac{\sqrt{1+x}-1}{x^2}$}
\prob{295}{$\displaystyle\lim_{x\to0}\frac{\sqrt{x^2+1}-1}{\sqrt{x^2+16}-4}$}
\prob{297}{$\displaystyle\lim_{x\to1}\frac{x^2-\sqrt x}{\sqrt x-1}$}
\prob{298}{$\displaystyle\lim_{h\to0}\frac{\sqrt{x+h}-\sqrt x}{h}$}
\prob{299}{$\displaystyle\lim_{x\to0}\frac{\sqrt[3]{1+x^2}-1}{x^2}$}
\prob{300}{$\displaystyle\lim_{x\to0}\frac{\sqrt[3]{1+x}-\sqrt[3]{1-x}}{x}$}
\prob{301}{$\displaystyle\lim_{x\to a}\frac{\sqrt{x-b}-\sqrt{a-b}}{x^2-a^2}$ \ ($a>b$)}
\prob{302}{$\displaystyle\lim_{x\to1}\frac{\sqrt[n]{x}-1}{\sqrt[m]{x}-1}$ \ ($n,m$ бүхэл)}
\prob{306}{$\displaystyle\lim_{x\to\infty}\bigl(\sqrt{x+a}-\sqrt x\bigr)$}
\prob{307}{$\displaystyle\lim_{x\to\infty}\bigl(\sqrt{x^2+1}-\sqrt{x^2-1}\bigr)$}
\prob{309}{$\displaystyle\lim_{x\to\pm\infty}x\bigl(\sqrt{x^2+1}-x\bigr)$ \ ($+\infty$ ба $-\infty$-г тусад нь)}
\end{plist}
\end{frame}
